\documentclass[sigconf,anonymous,nonacm]{aamas}

\usepackage{balance}
\usepackage{mathtools}
\usepackage{algorithm}
\usepackage[noend]{algpseudocode}
\usepackage{cleveref}
\usepackage{amsthm}

\settopmatter{printacmref=false}
\acmConference[]{}{}{}{}

\fancypagestyle{supplementstyle}{%
  \fancyhf{}
  \renewcommand{\headrulewidth}{0pt}
  \renewcommand{\footrulewidth}{0pt}
  \fancyfoot[C]{\thepage}
}

% Unnumbered restatements from the main paper
\newtheorem*{lemmafourone}{Lemma 4.1 (Main paper)}
\newtheorem*{theoremfourtwo}{Theorem 4.2 (Main paper)}

\newcommand{\alloc}{\mathcal{A}}
\newcommand{\allocB}{\mathcal{B}}

\title{Supplementary Material for
``Envy-Free Chore Allocation with Unit Subsidies under Negative
Dichotomous Valuations"}

\author{Anonymous Author(s)}

\begin{document}
\maketitle

\pagestyle{supplementstyle}
\thispagestyle{supplementstyle}




\section{Failure of an Incremental Unit-Subsidy Extension}
\label{sec:whynotbkns}


\begin{example}
\label{app:extension-failure}

Consider three agents $N=\{1,2,3\}$ and three chores $M=\{x,y,z\}$. We work in cost form, writing
\( c_i=-v_i \)
for agent \(i\)'s non-negative cost function. Thus, an allocation \(\mathcal A=(A_1,\ldots,A_n)\) together with a subsidy vector \(p=(p_1,\ldots,p_n)\) is envy-free iff
\[ c_i(A_i)-p_i\le c_i(A_j)-p_j
\qquad\text{for all } i,j\in N.
\]
The cost functions depend only on bundle size: \[ c_1(S)=\max\{0,|S|-1\}, \qquad c_2(S)=c_3(S)=\min\{|S|,2\}. \] Thus, \[
\begin{array}{c|cccc}
|S| & 0 & 1 & 2 & 3\\
\hline
c_1(S) & 0 & 0 & 1 & 2\\
c_2(S)=c_3(S) & 0 & 1 & 2 & 2 .
\end{array}
\] Every marginal cost belongs to $\{0,1\}$, so these are dichotomous cost functions. Suppose  $z$ is currently unallocated, and the partial allocation is \[
X_1=\{x,y\},
\qquad
X_2=X_3=\emptyset.
\] Note that this partial allocation is itself reachable from the empty allocation through successive insertions while maintaining the unit-subsidy invariant:
\[
(\emptyset,\emptyset,\emptyset)
\longrightarrow
(\{x\},\emptyset,\emptyset)
\longrightarrow
(\{x,y\},\emptyset,\emptyset),
\] with pointwise-minimal subsidy vectors
\(
(0,0,0), (0,0,0),\text{ and } (1,0,0),
\) respectively. The pointwise-minimal subsidy vector of $(X_1,X_2,X_3)$ is \[
p^*=(1,0,0).
\]

Indeed, agent $1$ has cost $1$ for her own bundle and cost $0$ for either empty bundle, so she requires one unit of subsidy. Agents $2$ and $3$ have no envy, since they have cost $0$ for their own empty bundles and cost $2$ for $X_1$.

We claim that the remaining chore $z$ cannot be inserted so that the resulting allocation admits an envy-free subsidy vector in $\{0,1\}^3$, even if the existing bundles $\{x,y\},\emptyset,\emptyset$ may first be permuted among the agents.

If $z$ is added to the two-chore bundle, the resulting bundle sizes are $(3,0,0)$. Every agent evaluates a three-chore bundle at cost $2$ and an empty bundle at cost $0$. Hence, the agent receiving the three-chore bundle requires a subsidy at least two units larger than that of an agent receiving an empty bundle. Thus, no subsidy vector in $\{0,1\}^3$ can make the resulting allocation envy-free.

It remains to consider the case in which $z$ is added to an empty bundle. The resulting bundle sizes are $(2,1,0)$, up to permutation. If an agent in $\{2,3\}$ receives the two-chore bundle, she evaluates it at cost $2$ and the empty bundle at cost $0$, again requiring a subsidy difference of at least two.

Therefore, the only remaining possibility is that agent $1$ receives the two-chore bundle. Let $i\in\{2,3\}$ receive the singleton and let $j$ receive the empty bundle. Envy-freeness for agent $1$ with respect to $i$ requires
\[
1-p_1\le 0-p_i,
\]
and hence
\[
p_1\ge p_i+1.
\]
Envy-freeness for agent $i$ with respect to $j$ requires \[
1-p_i\le -p_j,
\] and hence
\[
p_i\ge p_j+1.
\] Combining the two inequalities gives
\[
p_1\ge p_j+2,
\] which is impossible for $p\in\{0,1\}^3$. Hence, no assignment of the additional chore $z$, even after an arbitrary permutation of the previously allocated bundles, preserves the unit-subsidy guarantee. Notice that this is an obstruction to the incremental construction, not to the instance itself. The complete allocation
\[
(\{x\},\{y\},\{z\})
\] is envy-free without subsidies: agent $1$ evaluates every singleton at cost $0$, while agents $2$ and $3$ evaluate every singleton at cost $1$. \hfill $\square$
\end{example}



\section{Proof of Lemma 4.1 of the main paper}
We first restate the lemma for convenience.
\begin{lemmafourone}
Let $M'$ be a set of allocated items and let $A=(A_1,\ldots,A_n)$ be a complete allocation of $M'$. Suppose that the valuations are integer-valued and that there exists $p\in\{0,1\}^n$ such that $(A,p)$ is envy-free. Let $g\notin M'$ satisfy
\[
v_i(S\cup\{g\})-v_i(S)\in\{0,1\}
\]
for every agent $i$ and every $S\subseteq M'$. Then there exist a complete allocation $A'$ of $M'\cup\{g\}$ and $p'\in\{0,1\}^n$ such that $(A',p')$ is envy-free. Moreover, $(A',p')$ can be computed in polynomial time given value-oracle access.
\end{lemmafourone}

\begin{proof}
The proof is based on the one-good extension argument of Barman et al.~\cite{BKNS22}. Since that argument is stated there for dichotomous goods instances, we restate its relevant steps in a self-contained way and verify that each of them remains valid when the already allocated bundles may contain chores.


For an allocation $Y=(Y_1,\ldots,Y_n)$, let $G_Y$ denote its weighted envy graph, with edge weights
\[
w_Y(i,j):=v_i(Y_j)-v_i(Y_i).
\]
We use two standard characterizations of envy-freeability
\cite{HS19,BKNS22}. First, $Y$ is envy-freeable if and only if it maximizes
\[
\sum_{i\in N}v_i(Y_{\sigma(i)})
\]
over all permutations $\sigma$ of its bundles. Equivalently, $G_Y$ contains no positive-weight directed cycle. Second, if $Y$ is envy-freeable, then its pointwise-minimal supporting subsidy for agent $i$ is the maximum weight of a directed path starting at $i$ in $G_Y$.

We may assume without loss of generality that $p$ is the pointwise-minimal subsidy vector supporting $A$. Indeed, starting from the subsidy vector in the statement, replace it by the pointwise-minimal one. This can only decrease the subsidies componentwise. Since the valuations are integer-valued, all edge weights in $G_A$, and therefore all pointwise-minimal subsidies, are integers. Consequently, the resulting vector still belongs to $\{0,1\}^n$.

We first prove two properties of adding $g$ that will be used throughout the proof.

\medskip
\noindent
\textbf{Claim 1.}
Let $Y$ be an envy-freeable allocation of $M'$, and suppose that for some agent $x$,
\[
v_x(Y_x\cup\{g\})-v_x(Y_x)=1.
\]
Then
\[
Z=(Y_1,\ldots,Y_x\cup\{g\},\ldots,Y_n)
\]
is envy-freeable.

\smallskip
\noindent
\emph{Proof of Claim 1.}
Write
\[
\operatorname{SW}(Y):=\sum_{i\in N}v_i(Y_i).
\]
Since the marginal value of $g$ for $x$ on $Y_x$ is one,
\[
\operatorname{SW}(Z)=\operatorname{SW}(Y)+1.
\]
Consider any permutation $\sigma$ of the bundles of $Z$. Let $a$ be the agent who receives the bundle containing $g$ under this permutation. Removing $g$ from that bundle leaves a permutation of the bundles of $Y$. Since $Y$ is envy-freeable, the welfare of this reassignment is at most
$\operatorname{SW}(Y)$. Moreover,
\[
v_a(S\cup\{g\})-v_a(S)\leq 1
\]
for every $S\subseteq M'$. Hence
\[
\operatorname{SW}(Z_\sigma)
\leq \operatorname{SW}(Y)+1
=\operatorname{SW}(Z).
\]
Thus $Z$ maximizes welfare over all permutations of its bundles and is therefore envy-freeable.
\hfill$\square$

\medskip
\noindent
\textbf{Claim 2.}
Let $Y$ be an allocation and let
\[
Z=(Y_1,\ldots,Y_x\cup\{g\},\ldots,Y_n).
\]
Then
\[
w_Z(i,j)=w_Y(i,j)
\qquad
\text{for all }i,j\neq x,
\]
while
\[
w_Z(x,j)\leq w_Y(x,j)
\qquad\text{and}\qquad
w_Z(i,x)\leq w_Y(i,x)+1.
\]
Consequently, for every simple directed path $P$,
\begin{equation}\label{S1}
    w_Z(P)\leq w_Y(P)+1.
\end{equation}

\smallskip
\noindent
\emph{Proof of Claim 2.}
Only the bundle of $x$ changes. Hence, edges not incident to $x$ are unchanged. For every $j\neq x$,
\[
\begin{aligned}
w_Z(x,j)
&=v_x(Y_j)-v_x(Y_x\cup\{g\})\\
&=w_Y(x,j)
-\bigl(v_x(Y_x\cup\{g\})-v_x(Y_x)\bigr)
\leq w_Y(x,j),
\end{aligned}
\]
because the marginal value of $g$ is nonnegative. Similarly, for
$i\neq x$,
\[
\begin{aligned}
w_Z(i,x)
&=v_i(Y_x\cup\{g\})-v_i(Y_i)\\
&=w_Y(i,x)
+\bigl(v_i(Y_x\cup\{g\})-v_i(Y_x)\bigr)
\leq w_Y(i,x)+1.
\end{aligned}
\]
A simple path enters $x$ at most once. Its weight can therefore increase by at most one, proving \eqref{S1}.
\hfill$\square$

\medskip
We now follow the extendable/non-extendable distinction of Barman et al.~\cite{BKNS22}. For a subsidy vector $q$, write
\[
M(q):=\{i\in N:q_i\geq q_j\text{ for every }j\in N\}.
\]
Call $(A,p)$ \emph{extendable with $g$} if there exists a permutation $\sigma$ such that, writing
\[
B_i:=A_{\sigma(i)}
\qquad\text{and}\qquad
q_i:=p_{\sigma(i)},
\]
the pair $(B,q)$ is envy-free and there is some
$\kappa\in M(q)$ satisfying
\begin{equation}\label{S2}
 v_\kappa(B_\kappa\cup\{g\})-v_\kappa(B_\kappa)=1.   
\end{equation}


We will also use the following elementary permutation property. If $(A,p)$ is envy-free and $B=A_\sigma$ is envy-freeable, then $(B,p_\sigma)$ is envy-free. To see this, the envy-freeness of
$(A,p)$ gives
\[
v_i(A_i)+p_i
\geq
v_i(A_{\sigma(i)})+p_{\sigma(i)}
\qquad\text{for every }i.
\]
Summing these inequalities and using the fact that both $A$ and $B$ are envy-freeable shows that equality must hold for every $i$. The envy-freeness inequalities for $(B,p_\sigma)$ then follow directly. See Lemma 3 in Barman et al.~\cite{BKNS22} for a detailed argument. 

\medskip
\noindent
\textbf{Case 1: $(A,p)$ is extendable with $g$.}

Let $\sigma$ and $\kappa$ satisfy the above conditions, and define
\[
C=(B_1,\ldots,B_\kappa\cup\{g\},\ldots,B_n).
\]
By Claim~1, $C$ is envy-freeable.

It remains to show that $C$ can be supported by subsidies in $\{0,1\}^n$. Since $q_\kappa$ is maximal and
$q\in\{0,1\}^n$, there are two cases.

If $q_\kappa=0$, then $q_i=0$ for all $i\in N$. Thus, every directed path in $G_B$ has weight at most zero. By Claim~2, every directed path in $G_C$ has weight at most one. Hence, the pointwise-minimal supporting subsidies for $C$ are at most one. They are nonnegative integers, so they belong to $\{0,1\}^n$.

Suppose next that $q_\kappa=1$. Define
\[
\widehat q_\kappa:=0,
\qquad
\widehat q_i:=q_i
\quad (i\neq\kappa).
\]
We claim that $(C,\widehat q)$ is envy-free. For $i,j\neq\kappa$, nothing changes from $(B,q)$. For agent
$\kappa$, using \eqref{S2},
\[
v_\kappa(C_\kappa)+\widehat q_\kappa
=
v_\kappa(B_\kappa)+1
=
v_\kappa(B_\kappa)+q_\kappa
\geq
v_\kappa(B_j)+q_j.
\]
Finally, for every $i\neq\kappa$,
\[
\begin{aligned}
v_i(C_i)+\widehat q_i
&=v_i(B_i)+q_i\\
&\geq v_i(B_\kappa)+q_\kappa\\
&=v_i(B_\kappa)+1\\
&\geq v_i(B_\kappa\cup\{g\})
 =v_i(C_\kappa)+\widehat q_\kappa,
\end{aligned}
\]
where the last inequality follows because the marginal value of $g$ is at most one. Thus $\widehat q\in\{0,1\}^n$ supports $C$.

Hence, the desired extension exists in the extendable case.

\medskip
\noindent
\textbf{Case 2: $(A,p)$ is not extendable with $g$.}

We use the \textsc{FindSink} procedure of Barman et al.~\cite{BKNS22}, and verify that its proof remains valid in the
present setting.

For $s\in N$, write
\[
X^s:=(A_1,\ldots,A_s\cup\{g\},\ldots,A_n).
\]
Start with an arbitrary $s_0\in M(p)$. Whenever $X^{s_\tau}$ is envy-freeable, let $\phi^\tau$ denote its pointwise-minimal supporting subsidy vector. If every component of $\phi^\tau$ is strictly smaller than $2$, terminate. Otherwise, choose an agent
$s_{\tau+1}$ with
\[
\phi^\tau_{s_{\tau+1}}\geq2
\]
and continue with $X^{s_{\tau+1}}$.

We first establish the following:
\begin{equation}\label{S3}
    s_\tau\in M(p)
\quad\text{and}\quad
X^{s_\tau}\text{ is envy-freeable}
\end{equation}
for every agent considered by the procedure.

For $\tau=0$, the first part follows from the choice of $s_0$. Suppose, towards a contradiction, that $X^{s_0}$ is not
envy-freeable. By the welfare characterization, there exists a permutation of its bundles whose welfare is strictly larger. Since all valuations are integer-valued, the improvement is at
least one.

Let $B$ be the corresponding permutation of the bundles of $A$, obtained by deleting $g$ from the permuted bundle containing it, and let $k$ be the agent who receives the bundle $A_{s_0}$ in $B$. Since the marginal value of $g$ is nonnegative,
\[
\operatorname{SW}(X^{s_0})\geq\operatorname{SW}(A).
\]
Therefore,
\begin{equation}\label{S4}
\operatorname{SW}(B)
+
\bigl(v_k(B_k\cup\{g\})-v_k(B_k)\bigr)
\geq
\operatorname{SW}(A)+1.
\end{equation}
On the other hand, $A$ is envy-freeable, and hence
\[
\operatorname{SW}(B)\leq\operatorname{SW}(A).
\]
The marginal term in \eqref{S4} is at most one. Consequently, both inequalities must be tight:
\[
\operatorname{SW}(B)=\operatorname{SW}(A)
\]
and
\[
v_k(B_k\cup\{g\})-v_k(B_k)=1.
\]
Thus $B$ is envy-freeable. By the permutation property above, if $q$ is obtained by permuting $p$ in the same way, then $(B,q)$ is envy-free. Moreover, $B_k=A_{s_0}$ and $s_0\in M(p)$, so $k\in M(q)$. Hence $(A,p)$ is extendable with $g$, a contradiction. Therefore $X^{s_0}$ is envy-freeable.

Now suppose \eqref{S3} holds for $s_\tau$ and the procedure selects $s_{\tau+1}$ because $\phi^\tau_{s_{\tau+1}}\geq2$. We claim that $s_{\tau+1}\in M(p)$. If not, then the maximal subsidy in $p$ is one and $p_{s_{\tau+1}}=0$. Since $p$ is pointwise-minimal, every path starting from $s_{\tau+1}$ in $G_A$ has weight at most zero. By Claim~2, every path starting from $s_{\tau+1}$ in $G_{X^{s_\tau}}$ has weight at most one. Hence $\phi^\tau_{s_{\tau+1}}\leq1$, a contradiction. Thus $s_{\tau+1}\in M(p)$. The same argument used for $s_0$ then shows that $X^{s_{\tau+1}}$ is envy-freeable. This proves \eqref{S3}.

It remains to show that the procedure terminates. Suppose otherwise that some agent is selected twice. Consider two consecutive occurrences of a repeated agent and relabel the corresponding segment
as
\[
s_0,s_1,\ldots,s_T,
\qquad
s_T=s_0,
\]
with $s_0,\ldots,s_{T-1}$ distinct.

For every $\tau=1,\ldots,T$, agent $s_\tau$ is selected because its pointwise-minimal subsidy for $X^{s_{\tau-1}}$ is at least two. Hence, there exists a simple directed path $Q_\tau$ starting
at $s_\tau$ in $G_{X^{s_{\tau-1}}}$ such that
\begin{equation}\label{S6}
w_{X^{s_{\tau-1}}}(Q_\tau)\geq2.
\end{equation}
The path $Q_\tau$ must contain $s_{\tau-1}$. Otherwise, none of its edge weights differs from the corresponding edge weights in $G_A$, whereas every path starting at $s_\tau$ in $G_A$ has a weight at most 1, contradicting \eqref{S6}.



Let $P_\tau$ be the initial segment of $Q_\tau$ from $s_\tau$ to $s_{\tau-1}$. By Claim~2,
\begin{equation}\label{S7}
 w_A(Q_\tau)
\geq
w_{X^{s_{\tau-1}}}(Q_\tau)-1
\geq1.   
\end{equation}
The part of $Q_\tau$ following $P_\tau$ is a path starting at $s_{\tau-1}$ in $G_A$, and therefore has weight at most $p_{s_{\tau-1}}\leq1$. It follows from \eqref{S7} that
\begin{equation}\label{S8}
w_A(P_\tau)\geq0.
\end{equation}

Now concatenate the paths $P_T,P_{T-1},\ldots,P_1$. Since $s_T=s_0$, this concatenation starts and ends at $s_0$. If some vertex is repeated in the concatenation, split the resulting closed sequence of edges at repeated vertices. In this way, it can be decomposed into directed cycles of $G_A$. By \eqref{S8}, the total weight of the concatenated sequence is nonnegative. On the other hand, since $A$ is envy-freeable, $G_A$ contains no positive-weight directed cycle. Therefore, every cycle arising in the above decomposition has weight at most zero, while their total weight is nonnegative. It follows that the total weight is zero, and consequently, every cycle in the decomposition has weight zero.

We now derive a contradiction to non-extendability. Since $s_1$ was selected after considering $X^{s_0}$, there is a path $Q_1$ satisfying \eqref{S6}. Its weight in $G_A$ is at most $p_{s_1}\leq1$, while Claim~2 says that its weight can increase by at most one after $g$ is added to $A_{s_0}$. Hence
\[
w_{X^{s_0}}(Q_1)=2,
\qquad
w_A(Q_1)=1,
\]
and the weight of $Q_1$ increases by exactly one. Only an edge entering $s_0$ can increase when $g$ is added to $A_{s_0}$; edges leaving $s_0$ can only decrease. Therefore, if $k$ is the predecessor of $s_0$ on $Q_1$, then
\begin{equation}\label{S9}
v_k(A_{s_0}\cup\{g\})-v_k(A_{s_0})=1.
\end{equation}
Since the edge $(k,s_0)$ belongs to $P_1$, it belongs to one of the zero-weight directed cycles obtained above. Let $C$ be such a cycle. Cyclically permuting the bundles of $A$ along $C$ preserves social welfare, because the change in welfare is exactly $w_A(C)=0$. Let $\sigma$ denote the resulting permutation and let $B=A_\sigma$. Hence
\[
\operatorname{SW}(B)=\operatorname{SW}(A),
\] so $B$ is envy-freeable. By the permutation property, $(B,p_\sigma)$ is envy-free.

Furthermore, the cycle contains $(k,s_0)$, so agent $k$ receives bundle $A_{s_0}$ in $B$. Thus \eqref{S9} gives
\[
v_k(B_k\cup\{g\})-v_k(B_k)=1.
\]
Since $s_0\in M(p)$, agent $k$ belongs to $M(p_\sigma)$. Consequently, $\sigma$ and $k$ witness that $(A,p)$ is extendable with $g$, contradicting the assumption of Case~2.

Therefore, no agent is selected twice, and the procedure terminates after at most $n$ selections. At termination, the current allocation $X^s$ is envy-freeable by \eqref{S3}, and every component of its pointwise-minimal subsidy vector is strictly smaller than $2$. These subsidies are nonnegative integers, and hence they belong to $\{0,1\}^n$. This proves the existence of the desired allocation $A'$ and subsidy vector $p'$ in the non-extendable case.

Finally, all steps are polynomial-time. Pointwise-minimal subsidies can be computed from maximum-weight paths in an envy-freeable allocation. Extendability can be tested by the \textsc{Extend} subroutine of Barman et al.~\cite{BKNS22}: for each candidate recipient and most-subsidized bundle, it solves a maximum-weight bipartite matching problem. Its correctness uses only the welfare characterization of envy-freeability and the values of the existing bundles, and is therefore unchanged here. In the non-extendable case, the procedure above considers at most $n$ agents. Thus, the entire one-good extension can be carried out in polynomial time using value-oracle access.
\end{proof}

\section{Algorithm for Objective Signed Dichotomous Valuations}
We first restate the theorem for convenience.
\begin{theoremfourtwo}
For every instance with objective signed dichotomous valuations, there exist a complete allocation $\alloc$ and a subsidy vector $p\in\{0,1\}^n$ such that $(\alloc,p)$ is envy-free and
\[
\sum_{i\in N} p_i\leq n-1.
\]
Moreover, such an allocation and subsidy vector can be computed in
polynomial time in the value-oracle model.
\end{theoremfourtwo}
We next give the complete algorithm underlying Theorem 4.2 of the main paper, using Lemma 4.1 of the main paper as the one-good extension subroutine.

For brevity, let \textsc{ChoreAllocate} denote the polynomial-time procedure developed in Section~3 of the main paper for negative dichotomous valuations. Given a negative dichotomous chore instance, it returns a complete allocation $\alloc$ and $p\in\{0,1\}^n$ such that $(\alloc,p)$ is envy-free and $\sum_i p_i\le n-1$.

\begin{algorithm}[H]
\caption{Unit subsidies for objective signed dichotomous valuations}
\label{alg:objective-mixed}
\begin{algorithmic}[1]
\Require Agents $N=[n]$, items $M=G\mathbin{\dot\cup}C$, and value-oracle access to objective signed dichotomous valuations $\{v_i\}_{i\in N}$.
\Ensure A complete allocation $\alloc$ and $p\in\{0,1\}^n$ such that $(\alloc,p)$ is envy-free and $\sum_{i\in N}p_i\le n-1$.

\State Restrict each $v_i$ to subsets of $C$.
\State $(\alloc,p)\gets\textsc{ChoreAllocate}(N,C,\{v_i|_C\}_{i\in N})$.
\State Replace $p$ by the pointwise-minimal subsidy vector supporting $\alloc$.

\For{each $g\in G$}
    \State Apply the one-good extension procedure of  Lemma 4.1 of the main paper to $(\alloc,p)$ and $g$, obtaining an allocation $\alloc'$ of the currently allocated items together with $g$ and a subsidy vector $p'\in\{0,1\}^n$ such that $(\alloc',p')$ is envy-free.
    \State $\alloc\gets\alloc'$.
    \State Replace $p'$ by the pointwise-minimal subsidy vector supporting $\alloc$, and set $p\gets p'$.
\EndFor

\State \Return $(\alloc,p)$.
\end{algorithmic}
\end{algorithm}

The one-good extension in Line~5 is implemented by the \textsc{Extend}/\textsc{FindSink} procedure of Barman et al.~\cite{BKNS22}, whose validity in the present setting is established in the proof of Lemma 4.1 of the main paper. After every iteration, we take the pointwise-minimal supporting subsidy vector. Since the valuations are integer-valued, this vector is integral, and Lemma 4.1 of the main paper implies that all its components are at most one. Moreover, pointwise minimality implies that at least one component is zero. Hence, throughout the algorithm
\[
p\in\{0,1\}^n
\qquad\text{and}\qquad
\sum_{i\in N}p_i\le n-1.
\]

\section{Pareto optimality and unit subsidy bound for additive objective signed dichotomous valuations}
\begin{proposition}
\label{prop:additive-objective-po-supp}
For every instance with additive objective signed dichotomous valuations, there exist a complete allocation $\alloc$ and a subsidy vector $p\in\{0,1\}^n$ such that $(\alloc,p)$ is envy-free,
\[
\sum_{i\in N}p_i\leq n-1,
\]
and $\alloc$ is Pareto optimal. Moreover, such an outcome can be computed in polynomial time.
\end{proposition}

\begin{proof}
The proof has three steps. First, allocate the goods using Halpern and Shah's binary-additive result, obtaining an envy-freeable allocation with unit subsidies. Next, assign every chore that is costless to some agent to such an agent. Finally, assign each universally costly chore to an agent currently receiving zero subsidy and raise that agent's subsidy from $0$ to $1$; if all subsidies become $1$, subtract one from every subsidy. This preserves envy-freeness and keeps total subsidy at most $n-1$. Since every item is assigned to an agent attaining its maximum possible value, the final allocation is Pareto optimal.


We now formally write the details. Let $M=G\mathbin{\dot\cup}C$ be the common partition into goods and chores. Since the valuations are additive, for every agent $i$,
\[
v_i(g)\in\{0,1\}\qquad (g\in G)
\]
and
\[
v_i(c)\in\{0,-1\}\qquad (c\in C).
\]

We first allocate the goods. Ignore, for the moment, goods that have value zero for every agent, and apply the binary-additive goods result of Halpern and Shah~\cite{HS19} to the remaining goods. Their construction returns a non-wasteful allocation $\alloc^G$ that is envy-freeable and whose minimal subsidies have total at most $n-1$. Moreover, their proof shows that every path in the corresponding weighted envy graph has weight at most one. Since all edge weights are integral, the pointwise-minimal subsidies are integral as well, and hence we may choose
\[
p\in\{0,1\}^n,
\qquad
\sum_{i\in N}p_i\leq n-1.
\]
Goods having value zero for every agent may now be allocated arbitrarily, without affecting either envy-freeness or the subsidy vector.

We next allocate the chores. Partition $C$ into
\[
C^0:=\{c\in C:\text{there exists }i\in N\text{ with }v_i(c)=0\}
\]
and
\[
C^-:=\{c\in C:v_i(c)=-1\text{ for every }i\in N\}.
\]
Thus, $C^-$ is the set of universally costly chores.

First, consider a chore $c\in C^0$. Choose an agent $k$ with $v_k(c)=0$ and assign $c$ to $k$. This preserves envy-freeness with the same subsidy vector. Indeed, agent $k$'s value for her own bundle is unchanged, whereas for every agent $i\neq k$,
\[
v_i(A_k\cup\{c\})
=
v_i(A_k)+v_i(c)
\leq v_i(A_k).
\]
Hence, $k$'s bundle becomes weakly less desirable to every other agent, while all other relevant bundle values remain unchanged. Repeating this operation allocates all chores in $C^0$ while preserving the envy-free solution $(\alloc,p)$.

It remains to allocate the universally costly chores in $C^-$. We maintain the following invariant:
\begin{equation}\label{eq:additive-mixed-invariant}
(\alloc,p)\text{ is envy-free},\qquad
p\in\{0,1\}^n,
\qquad
\text{and }p_k=0\text{ for some }k\in N.
\end{equation}
The invariant holds initially. In particular, the pointwise-minimal subsidy vector has at least one zero component.

Let $c\in C^-$ be an unallocated chore, and choose an agent $k$ with
$p_k=0$. Assign $c$ to $k$ and define
\[
\widehat p_k:=1,
\qquad
\widehat p_i:=p_i\quad(i\neq k).
\]
We claim that the resulting solution remains envy-free.

Since $c$ has value $-1$ for every agent,
\[
v_k(A_k\cup\{c\})+\widehat p_k
=
v_k(A_k)-1+1
=
v_k(A_k)+p_k.
\]
Thus, agent $k$'s subsidized value for her own bundle is unchanged. Moreover, for every $i\neq k$,
\[
v_i(A_k\cup\{c\})+\widehat p_k
=
v_i(A_k)-1+1
=
v_i(A_k)+p_k.
\]
Hence, the subsidized value of $k$'s bundle is unchanged from the viewpoint of every agent. All other bundles and subsidies are unchanged. Consequently, every envy-freeness inequality is exactly the same as before, and the new solution is envy-free.

If $\widehat p$ has a zero component, set $p:=\widehat p$ and continue. Otherwise, if $\widehat p_i=1$ for all $i\in N$, replace it with
\[
p_i:=\widehat p_i-1=0 \text{ for all } i\in N.
\]
Subtracting the same amount from every agent's subsidy leaves all envy-freeness inequalities unchanged. Thus, the invariant \eqref{eq:additive-mixed-invariant} is restored. Repeating this procedure allocates every chore in $C^-$ and produces a complete allocation $\alloc$ with
\[
p\in\{0,1\}^n
\qquad\text{and}\qquad
\sum_{i\in N}p_i\leq n-1.
\]

It remains to prove Pareto optimality. By construction, every good that is valued at one by some agent is assigned to an agent who values it at one; a good valued at zero by every agent may be assigned arbitrarily. Every chore in $C^0$ is assigned to an agent who values it at zero, while every chore in $C^-$ has value $-1$ for every agent regardless of its recipient. Therefore, every item is assigned to an agent attaining the maximum possible value for that item. By additivity,
\[
\sum_{i\in N}v_i(A_i)
=
\sum_{e\in M}\max_{i\in N}v_i(e),
\]
which is the maximum possible utilitarian welfare over all complete allocations. Hence $\alloc$ is Pareto optimal.

Finally, the binary-additive goods allocation and its supporting subsidies can be computed in polynomial time by the construction of Halpern and Shah~\cite{HS19}. The subsequent chore-allocation steps require only identifying a zero-valued recipient for each chore and updating a binary subsidy vector. Hence, the entire construction runs in polynomial time.
\end{proof}

\bibliographystyle{ACM-Reference-Format}
\bibliography{references}


\end{document}
